\section{Related Work}
\label{sec:related}

\paragraph{Computer-use agent benchmarks.}
Benchmarks for agents that operate real software span three environment
families. On the desktop, OSWorld provides 369 tasks in executable virtual
machines with per-task evaluation scripts \citep{xie2024osworld}. On the web,
WebArena self-hosts realistic replicas of common sites \citep{zhou2024webarena},
while ClawBench evaluates agents directly on 144 live production websites,
capturing and blocking each task's final submission request so that real-world
side effects cannot occur \citep{zhang2026clawbench}. In the terminal,
Terminal-Bench curates hard, human-verified tasks with outcome-checking tests
\citep{merrill2026terminalbench}. Related environments cover mobile devices
\citep{rawles2024androidworld}, simulated workplaces \citep{xu2024agentcompany},
and repository-level software engineering \citep{jimenez2024swebench}.
Evaluations on live websites have proven harder than sandbox results suggest:
under human evaluation on real sites, reported web-agent success rates fall by
tens of points \citep{xue2025illusion}. All of these benchmarks evaluate a
single language, almost always English. AraClawBench borrows tasks from three of
them \citep{xie2024osworld, zhang2026clawbench, merrill2026terminalbench}
precisely so that its English arm stays comparable to established baselines
while its Arabic arm measures something no upstream benchmark covers.

\paragraph{Multilingual agent evaluation.}
A recent line of work extends agent benchmarks beyond English.
macOSWorld localizes an interactive macOS environment into five languages and
finds Arabic worst, with left-to-right click biases in mirrored interfaces
\citep{yang2025macosworld}. OmnilingualGAIA2 machine-translates a simulated
tool-use benchmark into ten languages and finds a universal, scale-resistant
gap concentrated in tool orchestration \citep{caciolai2026omnilingual}.
X-WebAgentBench translates a simulated shopping environment into fourteen
languages \citep{wang2025xwebagentbench}, MPR-GUI builds aligned multilingual
GUI perception tests \citep{chen2025mprgui}, MAPS extends four agentic
benchmarks into twelve languages and ties degradation to the share of
non-English input tokens \citep{hofman2025maps}, and PolyWorkBench authors
multilingual workplace workflows natively \citep{li2026polyworkbench}.
Two properties are shared by all of this work. First, with the partial
exception of PolyWorkBench, the non-English arms are produced by machine
translation; GAIA-v2-LILT demonstrates that correcting such arms with bilingual
human audit moves agent scores by up to 32.7 points, so translation quality is
load-bearing, not cosmetic \citep{kim2026lilt}. Second, every environment is a
sandbox or simulation; none observes an agent in a localized operating system
on live regional websites. No prior work evaluates any Arabic dialect, and none
includes tasks native to the target region rather than translated from Western
originals.

\paragraph{Arabic and dialectal evaluation.}
Arabic NLP evaluation has established that models handle Modern Standard
Arabic far better than dialects: AraDiCE builds dialectal and cultural
benchmarks by machine translation with extensive human post-editing and finds
that models understand dialects better than they generate them
\citep{mousi2024aradice}. This line of work is confined to static text tasks.
AraClawBench extends dialect evaluation to interactive agents: the same
computer-use task, issued in Egyptian, Syrian, Qatari, or Sudanese Arabic, on
the same Arabic machine.

\paragraph{Benchmark validity.}
A parallel literature examines whether agent benchmarks measure what they
claim. The Agentic Benchmark Checklist formalizes task validity and outcome
validity and finds violations in most popular agentic benchmarks, including
web-drift breakage in OSWorld's browser tasks \citep{zhu2025abc}.
\citet{shao2026protocolvalidity} formulate protocol validity --- a score
supports a capability claim only if the protocol keeps that capability
necessary for success --- and document reward hacking across fifteen
benchmarks. Task-design guidelines from Terminal-Bench reviewing emphasize
outcome-checking tests and adversarial review \citep{bercovich2026guidelines}.
For cross-language measurement these concerns compound: a language gap can be
manufactured by translation defects, evaluation artifacts, or environment
faults as easily as by model limitations \citep{kim2026lilt,
caciolai2026omnilingual}. AraClawBench's design --- frozen instructions with
enforced literal parity, deterministic evidence-based grading, evaluation
errors excluded from scores, and per-divergence adjudication --- is our answer
to exactly this compounding.

\medskip
\noindent In sum: the cross-lingual agent gap is established, Arabic is its
consistent worst case, and the existing evidence rests on machine-translated
tasks in simulated environments with no dialect coverage and no native tasks.
AraClawBench occupies that composite gap.
